\subsection{Multiplicity} 
What if we have a power of a factor, as in this polynomial?
\[
(x-4)(x+2)^2(x-7)^3(x+1)
\]
We saw in \unref{UnPolyBasics} that the equation
\[
(x-4)(x+2)^2(x-7)^3(x+1)=0
\]
only has four solutions:
\begin{lpic}{}{2}
\end{lpic}
But this is definitely a different polynomial from
\[
(x-4)(x+2)(x-7)(x+1)
\]
We can think of the factors with exponents---and therefore the roots that go with them---as appearing ``more than once.''
\begin{definition}[Multiplicity]
Recall for a polynomial $P$ and a number $r$, if $r$ is a root of $P$, then $(x-r)$ is a factor of $P$. If $(x-r)^m$ is a factor of $P$ (and $m$ is the largest number this applies to), we say the root $r$ has \term{multiplicity} $m$.
\end{definition}
In our example
\[
(x-4)(x+2)^2(x-7)^3(x+1)
\]
\begin{itemize}
\item The root $4$ has multiplicity \makeblank[.5in]
\item The root $-2$ has multiplicity \makeblank[.5in]
\item The root $7$ has multiplicity \makeblank[.5in]
\item The root $-1$ has multiplicity \makeblank[.5in]
\end{itemize}
\begin{Example}
Write a polynomial with leading coefficient 2 whose roots are 0 with multiplicity 2, -1 with multiplicity 3, 4 with multiplicity 1, and 3 with multiplicity 2.
\begin{lpic}{}{6}
\regtextleft{1}{-6}{The polynomial is};
\end{lpic}
\end{Example}